\appendix

\section{Auxiliary results}\label{appendixa}

In this appendix, we provide a concise overview of the resurgence properties of the asymptotic expansion of the gamma function and the uniform asymptotic expansion of the incomplete gamma function. Our focus is on the asymptotic behaviour of the coefficients of these expansions, as they play a significant role in proving our main result.

We define the scaled gamma function $\Gamma^\ast(a)$ through the relation
\begin{equation}\label{eq24}
\Gamma^\ast (a) = \frac{\Gamma (a)}{\sqrt {2\pi } a^{a - 1/2} \e^{ - a} }
\end{equation}
valid for $|\arg a|<\pi$. It is a well-established fact that both the scaled gamma function and its reciprocal possess asymptotic expansions given by
\begin{equation}\label{eq25}
\Gamma^\ast (a) \sim \sum_{n = 0}^\infty ( - 1)^n \frac{\gamma _n }{a^n } ,\qquad \frac{1}{\Gamma^\ast (a)} \sim \sum_{n = 0}^\infty \frac{\gamma _n }{a^n }
\end{equation}
as $a\to \infty$ in the sector $|\arg a| \le \pi-\delta<\pi$ (see, e.g., \cite[Section~3.6]{Temme1996}). Here, the $\gamma_n$ are known as the Stirling coefficients, with the initial values being
\[
\gamma _0 = 1,\qquad\gamma _1 = - \frac{1}{12},\qquad\gamma _2 = \frac{1}{288},\qquad\gamma _3 = \frac{139}{51840},\qquad\gamma _4 = - \frac{571}{2488320}.
\]
Note that the asymptotic expansion of the reciprocal scaled gamma function involves the same coefficients as that of the scaled gamma function, but with different signs of the coefficients with odd index. For an explanation of this phenomenon, refer to, for instance, \cite[p.~63]{Temme1996}. The asymptotic behaviour of the Stirling coefficients is described by the inverse factorial series
\begin{equation}\label{eq18}
\gamma _n \sim -\frac{1}{\pi}\sum_{k= 0}^{\infty} \gamma_k \sin\left( \frac{n -k}{2}\pi \right)\frac{\Gamma (n - k)}{(2\pi )^{n - k}}
\end{equation}
as $n\to+\infty$. Notably, the coefficients within the inverse factorial series are once again the Stirling coefficients, highlighting the resurgence property of the gamma function. The expansion \eqref{eq18} was established by Boyd \cite[equation~(3.34)]{Boyd1994}, although it had previously appeared without proof in an earlier paper by Rosser \cite[equation~(97)]{Rosser1955}.

Moving on to the incomplete gamma function $Q(a,z)$, define
\[
\lambda = \frac{z}{a},\qquad \eta = (2(\lambda - 1 - \log \lambda ))^{\frac{1}{2}} ,
\]
where the branch of the square root is continuous and satisfies $\eta(\lambda)\sim \lambda-1$ as $\lambda\to 1$. Temme~\cite{Temme1979} established that the incomplete gamma function admits the uniform asymptotic expansion
\begin{equation}\label{eq23}
Q(a,z) \sim \frac{1}{2}\erfc\bigg( \eta \sqrt{\frac{a}{2}} \bigg) + \frac{1}{\sqrt{2\pi a} }\exp \left(- \eta ^2 \frac{a}{2}\right)\sum_{n = 0}^\infty \frac{c_n (\eta )}{a^n},
\end{equation}
as $a\to \infty$ in the sector $|\arg a| \le \pi-\delta<\pi$, uniformly with respect to $\lambda$ in the sector $|\arg \lambda| \le 2\pi-\delta<2\pi$. The coefficients $c_n (\eta )$ are complicated functions of their argument with removable singularities at $\eta=0$. The resurgence properties of this expansion were studied by Olde Daalhuis~\cite{OldeDaalhuis1998a}. He demonstrated that the coefficients $c_n (\eta)$ possess inverse factorial series of the following form as $n\to+\infty$, for $\eta \in \mathcal{N}$:
\begin{gather}\label{eq22}
\begin{split}
c_n (\eta ) \sim \; & \frac{( - 1)^{n + 1} }{2\sqrt {2\pi }}\Gamma \left( n + \frac{1}{2} \right)\left( \frac{1}{2}\eta ^2 + 2\pi \im \right)^{ - n - \frac{1}{2}} I_{1 + \frac{\eta ^2 }{4\pi \im}} \left( n + \frac{1}{2},\frac{1}{2} \right)
\\ & + \frac{( - 1)^{n + 1} }{2\sqrt {2\pi } }\Gamma \left( n + \frac{1}{2} \right)\left( \frac{1}{2}\eta ^2 - 2\pi \im \right)^{ - n - \frac{1}{2}} I_{1 - \frac{\eta ^2 }{4\pi \im}} \left( n + \frac{1}{2},\frac{1}{2} \right)
\\ & - \frac{1}{\pi}\sum_{k = 0}^\infty c_k (\eta )\sin \left( \frac{n - k}{2}\pi \right)\frac{\Gamma (n - k)}{(2\pi )^{n - k} },
\end{split}
\end{gather}
where
\[
\mathcal{N} = \left\{ \eta \in \mathbb{C} \right\}\setminus \big\{ \eta = \alpha + \im \beta \mid \alpha \beta = \pm 2\pi ,\; \alpha \le - \sqrt {2\pi}\big\}
\]
is the image of the sector $|\arg \lambda|<2\pi$ under the mapping $\eta(\lambda)$. In \eqref{eq22}, $I_x(a,b)$ denotes the incomplete beta function \cite[Section~8.17]{DLMF}. In the special case that $\eta=0$, the expansion \eqref{eq22} simplifies to the inverse factorial series
\begin{equation}\label{eq19}
c_n (0) \sim -\frac{\Gamma \left( n + \frac{1}{2} \right)}{(2\pi )^{n + 1}}\sin\left( \left( n + \frac{1}{2} \right)\frac{\pi}{2} \right) - \frac{1}{\pi}\sum_{k = 0}^\infty c_k (0)\sin \left( \frac{n - k}{2}\pi \right)\frac{\Gamma (n - k)}{(2\pi )^{n - k} }
\end{equation}
as $n\to+\infty$.

An alternative asymptotic expansion for $c_n (\eta)$ was established in the paper \cite{Dunster1998}, but it does not exhibit the resurgence property.

\section[The computation of the coefficients C\_n(tau)]{The computation of the coefficients $\boldsymbol{C_n(\tau)}$}\label{appendixb}

In this appendix, we present an efficient method for computing the coefficients $C_n(\tau)$. We~substitute the polynomial expansion
\begin{equation}\label{eq29}
C_n(\tau) = \sum_{k=0}^{3n+2} c_{n,k}\tau^k
\end{equation}
into \eqref{eq32}. As a result, this substitution yields the following recurrence relation for the coefficients~$c_{n,k}$:
\begin{equation}\label{eq30}
c_{n,k}=(k+2)c_{n,k+2}+(k+1) c_{n-1,k+1}-\frac{2k}{k+1}c_{n-1,k-1}+\frac{1}{k+1}c_{n-1,k-3}.
\end{equation}
Given that $C_0 (\tau ) = \frac{1}{3}\tau^2 -\frac{1}{3}$, we deduce that
\begin{equation}\label{eq31}
c_{n,3n+2}=\frac{1}{3^{n+1}(n+1)!},\qquad c_{n,3n+1}=0
\end{equation}
holds for $n=0$. Using the recurrence relation \eqref{eq30}, we can demonstrate that \eqref{eq31} remains valid for $n\geq1$. Initiating with the first equation in \eqref{eq31} and subsequently setting $k=3n, 3n-2, \allowbreak 3n-4, \dots$, enables the computation of non-zero coefficients in the expansion \eqref{eq29}. Moreover, we establish that $c_{n,k}=0$ for $k=3n+1, 3n-1, 3n-3, \dots$. The following Wolfram Mathematica~\cite{Mathematica} code implements the recursive method:
\[
\begin{aligned}
&\bm{c[0,2] = 1/3;} \\
&\bm{c[0,0] = -1/3;} \\
&\bm{c[0,4] = 0;} \\
&\bm{c[0,-2] = 0;} \\
&\bm{\textbf{NN} = 10;} \\
&\bm{\textbf{For}[n = 1, n \leq \textbf{NN}, n\textbf{++},} \\
&\quad \bm{c[n,3*n+2] = 1/3{}^{\wedge}(n+1)/(n+1)!;} \\
&\quad \bm{c[n,3*n+4] = 0;} \\
&\quad \bm{\textbf{For}[k = 3*n, k \geq 0, k-=2,} \\
&\quad \quad \bm{c[n,k] = (k+2)*c[n,k+2] + (1/(k+1))*c[n-1,k-3]} \\
&\quad \quad \quad \bm{- (2k/(k+1))*c[n-1,k-1] + (k+1)*c[n-1,k+1];];} \\
&\quad \bm{c[n,k] = 0;} \\
&\quad \bm{c[n,k-2] = 0;];} \\
&\bm{\textbf{Ccoefficient}[\textbf{n$_{-}$},\textbf{t$_{-}$}]\,\textbf{:=}\,\textbf{Sum}[c[n,3*n+2-2*k]*t{}^{\wedge}(3*n+2-2*k),}\\
& \hphantom{\textbf{Ccoefficient}[\textbf{n$_{-}$},\textbf{t$_{-}$}]\,\textbf{:=}\,} \bm{\{k,0,3/2*n+1\}];} \\
&\bm{\textbf{Table}[\textbf{Ccoefficient}[n,t],\{n,0,\textbf{NN}\}]}
\end{aligned}
\]
This code generates a list of coefficients $C_n(t)$ for $0\le n\le 10$. Adjusting the value of the variable \textbf{NN} allows for computing a different number of coefficients.
